\documentclass[../../../include/open-logic-section]{subfiles}

\begin{document}

\olfileid{sth}{cardinals}{hp}
\olsection{ملحق: مبدأ هيوم}

وصفنا في \olref[cp]{sec} مبدأ كانتور، وكان كما يأتي:
\begin{align*}
	\card{A} = \card{B} & \text{ إذا وفقط إذا كان } A \approx B.
\intertext{وهذا شبيه جداً بما يسمى اليوم
\emph{مبدأ هيوم}، الذي يقول:}
	\fregenum{x} {F(x)} = \fregenum{x}{G(x)} & \text{ إذا وفقط إذا كان } F \sim G
\end{align*}
حيث تختصر «$F \sim G$» القول إن عدد الـ$F$ات يساوي تماماً عدد
الـ$G$ات، أي إنه يمكن إقامة تقابل بين الـ$F$ات والـ$G$ات، أي:
\begin{align*}
	\exists R(&\forall v\forall y(Rvy \lif (Fv \land Gy)) \land {}\\
		&\forall v(Fv \lif \lexists![y][Rvy]) \land {}\\
		&\forall y(Gy \lif \lexists![v][Rvy]))
\end{align*}
لكن بين مبدأ هيوم ومبدأ كانتور فرقاً في النوع. ففي عبارة مبدأ كانتور،
المتغيران «$A$» و«$B$» حدّان من الرتبة الأولى يدلان على
\emph{مجموعتين}. أما في عبارة مبدأ هيوم، فإن «$F$» و«$G$» و«$R$»
\emph{ليست} حدوداً من الرتبة الأولى؛ بل تقع في \emph{موضع المحمول}.
(ولعلها تدل على \emph{خصائص}.) ولذلك يمكن أن نشرح مبدأ هيوم بالعربية
كما يأتي: عدد الـ$F$ات يساوي عدد الـ$G$ات إذا وفقط إذا أمكن إقامة
تقابل بين الـ$F$ات والـ$G$ات. وسمي هذا \emph{مبدأ هيوم} لأن هيوم كتب
ذات مرة:
\begin{quote}
  حين يُقرن عددان بحيث تقابل دائماً كل وحدة من أحدهما وحدة من الآخر،
  نحكم بأنهما متساويان.
  \citep[Pt.III Bk.1 \S1]{Hume1740}
\end{quote}
وقد منح \citet[\S63]{Frege1884} مبدأ هيوم مكانته البارزة في المنطق
الرياضي المعاصر؛ إذ اقتبس هذه الفقرة من هيوم قبل أن يرسم، في الواقع،
ملامح ما سميناه مبدأ هيوم.

ينبغي أن تلاحظ التشابه البنيوي بين مبدأ هيوم والقانون الأساسي الخامس.
وقد صغنا الأخير في \olref[story][blv]{sec} كما يأتي:
\[
	\fregeext{x}{F(x)} = \fregeext{x}{G(x)} \text{إذا وفقط إذا كان } \lforall[x][(F(x) \liff G(x))].
\]
وقد يفيد هنا بعض التعليق والمقارنة.

ثمة طريقتان لفهم مبدأ كمبدأ هيوم أو القانون الأساسي الخامس: فهماً
\emph{حملياً} أو فهماً \emph{غير حملي} (راجع
\olref[story][predicative]{sec}). وفق القراءة غير الحملية للقانون الأساسي
الخامس، يقع الكائن $\fregeext{x}{F(x)}$، لكل $F$، ضمن مجال التكميم الذي
استعملناه في صياغة القانون الأساسي الخامس نفسه. وبالمثل، وفق القراءة
غير الحملية لمبدأ هيوم، يقع الكائن $\fregenum{x}{F(x)}$، لكل $F$، ضمن
مجال التكميم الذي استعملناه في صياغة مبدأ هيوم. أما وفق الفهم
\emph{الحملي}، فيكون الكائنان $\fregeext{x}{F(x)}$
و$\fregenum{x}{F(x)}$ عنصرين من مجال \emph{مختلف}.

إذا قرأنا القانون الأساسي الخامس قراءة غير حملية، فإنه يفضي إلى عدم
الاتساق بواسطة الاستيعاب الساذج (للتفاصيل، انظر
\olref[story][blv]{sec}). ويمكن، شأنه شأن الاستيعاب الساذج، جعله متسقاً
بقراءته \emph{حملياً}. لكنه على الأرجح لن يحقق كل ما أردناه منه.

أما مبدأ هيوم، فيمكن \emph{بصورة متسقة} أن يقرأ قراءة غير حملية، وهو
بهذه القراءة قوي جداً.

وللتوضيح، انظر إلى المحمول «$x \neq x$»، الذي لا يحققه أي شيء بوضوح.
يعطينا مبدأ هيوم عندئذ كائناً هو $\# x( x\neq x)$. ويمكن أن نعامل هذا
على أنه العدد~$0$. والآن، وفق الفهم \emph{غير الحملي}---ولكن وفق هذا
الفهم \emph{وحده}---يقع هذا الكائن $0$ ضمن مجال التكميم الأصلي. لذلك
يمكننا أن نطبق مبدأ هيوم تطبيقاً ذا معنى على المحمول «$x = 0$» فنحصل
على كائن $\#x (x = 0)$. ويمكن أن نعامل هذا على أنه العدد~$1$. وفوق ذلك،
يستلزم مبدأ هيوم أن $0 \neq 1$، إذ لا يمكن إقامة تقابل من الكائنات غير
المطابقة لذواتها إلى الكائنات المطابقة لـ$0$ (لا يوجد أي من الأولى،
ويوجد واحد من الثانية). ثم إن $1$، إذا واصلنا العمل على نحو غير حملي،
يقع ضمن مجال التكميم الأصلي. فيمكننا أن نطبق مبدأ هيوم تطبيقاً ذا معنى
على المحمول «$(x = 0 \lor x = 1)$» فنحصل على كائن
$\#x(x = 0 \lor x = 1)$. ويمكن أن نعامل هذا على أنه العدد~$2$، ونستطيع
أن نبين أن $0\neq 2$ و$1 \neq 2$، وهكذا.

وخلاصة القول إن مبدأ هيوم، إذا أُخذ على نحو غير حملي، يستلزم وجود
\emph{عدد لا متناهٍ من الكائنات}. وقد شجع ذلك
\emph{المناطقة الفريغيين الجدد} على اتخاذ مبدأ هيوم أساساً للحساب.

غير أن فرغه \emph{نفسه} لم يتخذ مبدأ هيوم أساساً للحساب. بل أثبته من
تعريف صريح: يُعرَّف $\fregenum{x}{F(x)}$ بأنه امتداد المفهوم
$F \sim \Phi$. وبمصطلحات حديثة، قد نحاول صياغة ذلك هكذا:
$\fregenum{x}{F(x)} = \Setabs{G}{F \sim G}$؛ لكن هذا يعيدنا إلى مشكلات
الاستيعاب الساذج.

\end{document}
